\subsection{Application to Field Theory}

PROPOSITION 9. --- Let L be a field and E a subring of \(\operatorname{End}_{\mathbf{Z}}(\mathrm{L})\) that contains the mapping \(\gamma_{a}: x \mapsto ax\) for every \(a \in L\) . Denote by K the set of elements a of L such that \(u(xa) = u(x)a\) for every x in L and every u in E; it is a subfield of L.

Let V be a finite-dimensional linear subspace of the right K-vector space L, and let h be a K-linear mapping from V to L. There exists an element of E that coincides with h on V.

We view L as a left E-module. Since E contains the left multiplications \(\gamma_{a}\) , every E-submodule of L is a left ideal of the field L; the E-module L is therefore simple. Every endomorphism of the additive group of L that commutes with the \(\gamma_{a}\) is of the form \(\delta_{b}: x \mapsto xb\) with b in L. Consequently, \(b \mapsto \delta_{b}\) is an isomorphism from K to the opposite ring of \(\operatorname{End}_{\operatorname{E}}(\operatorname{L})\) , which is a field.

The bicommutant \(E''\) of the E-module L therefore consists of the endomorphisms of the right K-vector space L. Let v be an endomorphism of the K-vector space L whose restriction to V coincides with h; it is an element of \(E''\) . Let \((x_{i})_{i\in\mathrm{I}}\) be a basis of V over K; by Theorem 3 (VIII, p.~88), there exists an element u of E such that \(u(x_{i}) = v(x_{i}) = h(x_{i})\) for \(i \in I\) . By linearity, it follows that \(u(x) = h(x)\) for every x in V.

COROLLARY. --- Let L be a field. Let \(\Gamma\) be a subgroup of the automorphism group of the field L, and let K be the field of invariants of \(\Gamma\) . Let V be a right K-linear subspace of L of finite dimension n over K. Then there exist elements \(\sigma_{1},\ldots,\sigma_{n}\) of \(\Gamma\) with the following property: for every K-linear mapping u from V to L, there exist elements \(a_{1},\ldots,a_{n}\) of L such that we have \(u(x)=\sum_{i=1}^{n}a_{i}\sigma_{i}(x)\) for every x in V.

Denote by E the set of mappings from L to L of the form \(x \mapsto \Sigma_{\sigma \in \Gamma} a_{\sigma} \sigma(x)\) , where \((a_{\sigma})_{\sigma \in \Gamma}\) is a family of elements of L with finite support. We have \(\gamma_{a} \in E\) for every a in L, and E is a subring of the endomorphism ring of the additive group of L. Moreover, the field K consists of the elements a of L such that \(u(xa) = u(x)a\) for every x in L and every u in E.

Let H be the left L-vector space \(\operatorname{Hom}_{\mathrm{K}}(\mathrm{V},\mathrm{L})\) ; it has dimension n.~By Proposition 9, it is generated by the restrictions of the elements of \(\Gamma\) to V. There exist n elements \(\sigma_{1},\ldots,\sigma_{n}\) of \(\Gamma\) whose restrictions to V form a basis of H over L. The corollary follows from this.

Remark. --- When the field L is commutative, this corollary reduces to Artin's theorem (V, §10, No.~6, p.~65).

